\documentclass[10pt,a4paper]{article}
\usepackage[utf8]{inputenc}
\usepackage{amsmath}
\usepackage{amsfonts}
\usepackage{amssymb}
\usepackage{makeidx}
\usepackage{graphicx}
\usepackage[width=15.00cm, height=22.00cm]{geometry}
\author{José Ignacio Escribano}
\title{Algunas definiciones}

\usepackage{amsmath}
\usepackage{xparse}
\usepackage{amsthm}
\usepackage{tikz}

% Seno
\def\sen{\mathop{\mbox{\normalfont sen}}\nolimits}

% Seno hiperbólico
\def\senh{\mathop{\mbox{\normalfont senh}}\nolimits}

% Arcoseno
\def\arcsen{\mathop{\mbox{\normalfont arcsen}}\nolimits}

% Gradiente
\def\grad{\mathop{\mbox{\normalfont grad}}\nolimits}

% Argumento principal
\def\Arg{\mathop{\mbox{\normalfont Arg}}\nolimits}

% Logaritmo complejo
\def\Log{\mathop{\mbox{\normalfont Log}}\nolimits}

% Rotacional
\def\rot{\mathop{\mbox{\normalfont rot}}\nolimits}

% Divergencia
\def\diver{\mathop{\mbox{\normalfont div}}\nolimits}

% Longitud
\def\longi{\mathop{\mbox{\normalfont long}}\nolimits}

% Índice de una curva
\def\ind{\mathop{\mbox{\normalfont Ind}}\nolimits}

% Rango de una matriz
\def\rang{\mathop{\mbox{\normalfont rang}}\nolimits}

% Norma
\newcommand{\norm}[1]{\lVert #1 \rVert}

% Derivada parcial
\newcommand{\derpar}[1]{\dfrac{\partial}{\partial #1}}

% Derivada parcial con función
\newcommand{\derparcial}[2]{\dfrac{\partial #1}{\partial #2}}

% Derivada parcial de orden n con funcion
\newcommand{\derparcialn}[3]{\dfrac{\partial^{#3} #1}{\partial #2^{#3}}}

% Integral doble en [a,b] x [c,d]
\newcommand{\intdob}[4]{\displaystyle \int_{#1}^{#2} \int_{#3}^{#4}}

% Integral triple en [a,b] x [c,d] x [e,f]
\newcommand{\inttri}[6]{\displaystyle \int_{#1}^{#2} \int_{#3}^{#4} \int_{#5}^{#6}}

% Números complejos
\def\C{\ensuremath{\mathbb{C}}}

% Números reales
\def\R{\ensuremath{\mathbb{R}}}

% Números racionales
\def\Q{\ensuremath{\mathbb{Q}}}

% Números enteros
\def\Z{\ensuremath{\mathbb{Z}}}

% Números naturales
\def\N{\ensuremath{\mathbb{N}}}

% Teorema
\newtheorem{teo}{Teorema}

% Corolario
\newtheorem{cor}{Corolario}

% Proposición
\newtheorem{prop}{Proposición}

% Lema
\newtheorem{lema}{Lema}

% Definición
\newtheorem{defi}{Definición}

% Observaciones
\newtheorem*{obs}{Observaciones}

\newtheorem*{ob}{Observación}

\newtheorem*{nota}{Nota}

\newtheorem*{problema}{Problema}

\newtheorem*{pregunta}{Pregunta}

\newtheorem{ejemplo}{Ejemplo}

\NewDocumentCommand{\overarrow}{O{=} O{\uparrow} m}{%
  \overset{\makebox[0pt]{\begin{tabular}{@{}c@{}}#3\\[0pt]\ensuremath{#2}\end{tabular}}}{#1}
}
\NewDocumentCommand{\underarrow}{O{=} O{\downarrow} m}{%
  \underset{\makebox[0pt]{\begin{tabular}{@{}c@{}}\ensuremath{#2}\\[0pt]#3\end{tabular}}}{\ensuremath{#1}}
}

\newcommand{\contradiction}{%
\begin{tikzpicture}[rotate=45,x=0.5ex,y=0.5ex]
\draw[color=red, line width=.1ex] (0,2) -- (3,2) (0,1) -- (3,1) (1,3) -- (1,0) (2,3) -- (2,0);
\end{tikzpicture}
}


\begin{document}
\section{Grafos}

\begin{defi}
\textcolor{red}{Grafo}
\end{defi}

\begin{defi}
\textcolor{red}{Orden de un grafo}
\end{defi}

\begin{defi}
\textcolor{red}{Grado de un grafo}
\end{defi}

\begin{defi}
La fuerza de un nodo en un grafo ponderado es la suma de los pesos de los arcos incidentes.
\end{defi}

\section{Rankings}

\begin{defi}
Dado un conjunto $\mathcal{N} = \{1,\dots,n\}$ que llamaremos nodos, definimos un ranking $c$ de $\mathcal{N}$ como cualquier biyección $c: \mathcal{N} \to \mathcal{N}$.\\
Escribiremos $i \prec_c j$ cuando el nodo $i$ aparezca antes que el nodo $j$ en el ranking $c$. 
\end{defi}

\begin{defi}
Dada una familia finita $\mathcal{R} = \{c_1,c_2,\dots, c_r\}$ de rankings, decimos que el par de nodos $(i,j) \in \mathcal{N}$ ``compiten'' si existe $t \in \{1,2,\dots, r-1\}$ tal que $i$ y $j$ intercambian sus posiciones relativas entre dos rankings consecutivos $c_t$ y $c_{t+1}$.
\end{defi}

\begin{defi}
Definimos el grafo de competitividad de la familia de rankings $R$, denotado como $G_c(\mathcal{R}) = (\mathcal{N}, E_\mathcal{R})$, donde $E_\mathcal{R}$ denota el conjunto de arcos, como el grafo no dirigido con nodos $\mathcal{N}$ y arcos dados por la siguiente regla: hay un enlace entre $i$ y $j$ si $(i,j)$ ``compiten''.
\end{defi}

\begin{defi}
Decimos que dos nodos $i$, $j$ ``compiten'' $k$ veces, si $k$ es el número máximo de rankings donde $i$ y $j$ ``compiten''.
\end{defi}

\begin{defi}
Definimos el grafo de competitividad evolutivo de $\mathcal{R}$, denotado como $G_c^e(\mathcal{R}) = (\mathcal{N}, E_\mathcal{R}^e)$, como el grafo no dirigido ponderado con nodos $\mathcal{N}$ y arcos dados por la siguiente regla: hay un arco entre $i$ y $j$ etiquetado con peso $k$ si $(i,j)$ compiten $k$ veces.  
\end{defi}

\begin{nota}
Notar que la red no ponderada detrás del grafo ponderado $G_c^e(\mathcal{R})$ es $G_c(\mathcal{R})$.
\end{nota}

\begin{nota}
Notar que el orden de los rankings es fundamental en el cálculo de los pesos del grafo de competitividad evolutivo, aunque no tiene influencia en el grafo de competitividad (no ponderado). 
\end{nota}

\begin{defi}
Definimos el grado medio normalizado (normalized mean degree) de una familia de rankings $\mathcal{R}$ como la suma de todos los grados del grafo competitivo $G_c(\mathcal{R})$, dividido por la suma 
\end{defi}
\end{document}